\section{Step 69 Check the actual selected-level reconstruction}\label{h263-block-s09}
Run the public feedback on a source containing your block entry and completed
transform/quantizer dependencies:
\begin{verbatim}
python3 construction/h263-reconstruction-construction/grade.py \
    my-block.asm --cli /path/to/mars-cli.jar \
    --report block-results.json
\end{verbatim}
Compare every stage, not just the first final pixel. The checker accepts the
permitted forward and inverse accuracy while requiring exact level-to-coefficient
and inverse-to-pixel consistency. It also checks wide pointers, independent
strides, input-read bounds, output guards and saved state.

The statistical inverse gate remains a separate prerequisite. An implementation
that passes one constant block can still have a normalization, transpose or bias
error elsewhere.

\paragraph{Observe} Two independently accurate precisions and another traversal
order can satisfy the same block interface without sharing every rounded coefficient.
\Needspace{8\baselineskip}
\paragraph{Check} Why reconstruct the reference inverse from returned levels?
\begin{itemize}
  \item[A.] It forces every encoder to emit the host encoder's bits.
  \item[B.] It avoids checking quantization entirely.
  \item[C.] It checks the actual representation this block selected.
\end{itemize}
